\problemname{Inflation}
En stor del av ekonomin styrs av den rådande inflationen, d.v.s. hur snabbt priserna på varor i samhället stiger.
T.ex. påverkar hur stor inflationen är hur Riksbanken sätter sin styrränta, vilket i sin tur har en inverkan på aktiemarknaden, bolåneräntor, o.s.v.
Inflationen mäts i procent, och anger hur mycket priserna ökat jämfört med samma tidpunkt föregående år.

I Sverige är det myndigheten Statistiska centralbyrån (SCB) som mäter inflationstakten i samhället.
De beräknar varje månad hur stor inflationen var, genom att mäta hur inflationstakten \emph{ändras varje dag} under månaden och sedan slår samman alla dagsökningar.
Du vill gärna använda dessa siffror för att veta om du borde köpa eller sälja alla dina aktier innan SCB släpper siffrorna officiellt.
Som tur är har du en kompis som gärna läcker siffrorna till dig.

Ett problem är att din kompis varit rätt så inkonsekvent med sina anteckningar.
Ibland har han antecknat dagsändringen i inflationstakten i procentenheter och ibland i procent.
Om inflationstakten var $20\%$ och den antecknade ökningen var "$50$" vet du alltså inte om den nya inflationstakten är $70\%$ ($+50$ procentenheter) eller $30\%$ ($+50$ procent).
Dessutom har han skrivit sina anteckningar huller om buller, så du kan inte avgöra vilken ökning som hör till vilken dag i månaden.

Du kan åtminstone beräkna en övre gräns på inflationstakten, vilket kanske hjälper din aktiehandel något.
Vad är den högsta möjliga inflationstakten, givet din kompis data?

\section*{Indata}
Den första raden innehåller ett decimaltal $I$, inflationstakten i procent i början av månaden.
Detta följs av en rad med $30$ nya decimaltal, ändringarna i inflationstakten som din kompis antecknat under månaden.

Alla tal är mellan $0$ och $1000$ och har högst $2$ decimaler.


\section*{Utdata}
Skriv ut ett decimaltal: den högsta möjliga inflationstakten.

\section*{Poängsättning}
Din lösning kommer att testas på en mängd testfallsgrupper.
För att få poäng för en grupp så måste du klara alla testfall i gruppen.

\noindent
\begin{tabular}{| l | l | p{12cm} |}
  \hline
  \textbf{Grupp} & \textbf{Poäng} & \textbf{Gränser} \\ \hline
  $1$    & $22$      & Din kompis har antecknat alla ökningar i procentenheter. \\ \hline
  $2$    & $33$      & Din kompis har antecknat ökningarna i rätt ordning. \\ \hline
  $3$    & $45$      & Inga ytterligare begräsningar. \\ \hline
\end{tabular}

\section*{Förklaring av exempelfall}
TODO
